\section{Evaluation Results and Model Comparison}
\label{sec:evaluation}

\subsection{Test Set Evaluation}

All three models were evaluated on a held-out test set of 103 files from 19
patients (patient-wise split, no data leakage). Two threshold regimes are
reported: \textbf{optimal thresholds} (sweep 0.05--0.95, step 0.01, maximising
F1) and \textbf{fixed threshold 0.5} (the standard operating point).

\subsubsection{Baseline GAT (OptimizedRespiratoryGAT)}

The Baseline GAT was evaluated at both threshold regimes. At the standard
$t=0.5$ threshold:

\begin{table}[!h]
  \renewcommand{\arraystretch}{1.15}
  \caption{Baseline GAT Test Performance (threshold = 0.5)}
  \label{tab:baseline_gat_fixed}
  \centering
  \begin{tabular}{@{}lcc@{}}
    \toprule
    \textbf{Metric} & \textbf{Wheeze} & \textbf{Crackle} \\
    \midrule
    Precision        & 0.389 & 0.509 \\
    Recall           & 0.345 & 0.671 \\
    Specificity      & 0.757 & 0.519 \\
    F1-Score         & 0.366 & 0.579 \\
    ROC-AUC          & 0.611 & 0.634 \\
    PR-AUC           & 0.363 & 0.551 \\
    \bottomrule
  \end{tabular}
\end{table}

A constrained dual-threshold optimisation was conducted, finding task-specific
thresholds that meet clinical constraints (min precision = 0.35 wheeze, 0.45
crackle; min recall = 0.60 wheeze, 0.55 crackle). The optimisation achieved a
0.52 threshold for crackle meeting all constraints, F1 = 0.507 with 0.454
precision and 0.574 recall. However, for wheeze no thresholds met all
constraints; the fallback threshold of 0.28 provided a recall of 0.953 at a
precision of 0.346.

\paragraph{Best chunk-level performance at $t=0.5$.}
The Baseline GAT achieved the highest chunk-level F1 at the standard threshold:
Wheeze F1 = 0.544, Crackle F1 = 0.731 (Table~\ref{tab:baseline_gat_fixed}),
outperforming both GraphSAGE and Improved GAT at this operating point.

\subsubsection{GraphSAGE}

The GraphSAGE model was evaluated at both frame-level (0.5-second) and
chunk-level (5-second windows). The chunk-level decision rule required at
least 2 consecutive positive frames within a 5-second window.

\begin{table}[!h]
  \renewcommand{\arraystretch}{1.15}
  \caption{GraphSAGE Test Performance (Frame-Level)}
  \label{tab:graphsage_frame}
  \centering
  \begin{tabular}{@{}lcc@{}}
    \toprule
    \textbf{Metric} & \textbf{Wheeze} & \textbf{Crackle} \\
    \midrule
    Precision        & 0.389 & 0.556 \\
    Recall           & 0.460 & 0.503 \\
    Specificity      & 0.665 & 0.699 \\
    F1-Score         & 0.421 & 0.529 \\
    ROC-AUC          & 0.608 & 0.643 \\
    PR-AUC           & 0.367 & 0.533 \\
    \bottomrule
  \end{tabular}
\end{table}

\begin{table}[!h]
  \renewcommand{\arraystretch}{1.15}
  \caption{GraphSAGE Test Performance (Chunk-Level, 5s windows)}
  \label{tab:graphsage_chunk}
  \centering
  \begin{tabular}{@{}lcc@{}}
    \toprule
    \textbf{Metric} & \textbf{Wheeze} & \textbf{Crackle} \\
    \midrule
    Precision        & 0.386 & 0.647 \\
    Recall           & 0.611 & 0.611 \\
    Specificity      & 0.478 & 0.633 \\
    F1-Score         & 0.473 & 0.629 \\
    ROC-AUC          & 0.544 & 0.622 \\
    \bottomrule
  \end{tabular}
\end{table}

GraphSAGE achieved the \textbf{highest specificity at $t=0.5$}: 0.665 for
wheeze and 0.699 for crackle, demonstrating the best precision--recall
trade-off among all models. Temperature scaling ($T_w = 0.504$, $T_c = 2.977$)
ensured calibrated probability outputs.

\subsubsection{ImprovedRespiratoryGAT}

\begin{table}[!h]
  \renewcommand{\arraystretch}{1.15}
  \caption{Improved GAT Test Performance (threshold = 0.5)}
  \label{tab:improved_gat_05}
  \centering
  \begin{tabular}{@{}lcc@{}}
    \toprule
    \textbf{Metric} & \textbf{Wheeze} & \textbf{Crackle} \\
    \midrule
    Precision        & 0.000 & 0.000 \\
    Recall           & 0.000 & 0.000 \\
    Specificity      & 1.000 & 1.000 \\
    F1-Score         & 0.000 & 0.000 \\
    ROC-AUC          & 0.590 & 0.364 \\
    PR-AUC           & 0.478 & 0.539 \\
    \bottomrule
  \end{tabular}
\end{table}

At the standard $t=0.5$ threshold, the Improved GAT produces
\textbf{zero positive predictions}---all predicted probabilities fall below
0.5, rendering the model unable to detect any events. At its optimal threshold
($t=0.05$), it achieves 100\% recall at the cost of 0\% specificity. The
inverted crackle ROC-AUC (0.364) indicates unreliable probability ordering.

\paragraph{Quantitative Comparison}
Table~\ref{tab:model_comparison} summarises performance metrics and
architectural features for all three models.

\begin{table}[!t]
  \renewcommand{\arraystretch}{1.25}
  \caption{Comparative Evaluation of Trained GNN Architectures}
  \label{tab:model_comparison}
  \centering
  \begin{tabular}{@{}lccc@{}}
    \toprule
    \textbf{Metric / Feature}
      & \textbf{Optimized}
      & \textbf{GraphSAGE}
      & \textbf{Improved} \\
    \midrule
    Val.\ Weighted F1          & 0.7667      & ---         & \textbf{0.9384} \\
    Best Training Loss         & $\sim$0.15  & ---         & \textbf{0.0679} \\
    \midrule
    Frame Wheeze F1 ($t=0.5$)  & 0.366       & \textbf{0.421} & 0.000 \\
    Frame Crackle F1 ($t=0.5$) & \textbf{0.579} & 0.529     & 0.000 \\
    Chunk Wheeze F1 ($t=0.5$)  & \textbf{0.544} & 0.473     & 0.000 \\
    Chunk Crackle F1 ($t=0.5$) & \textbf{0.731} & 0.629     & 0.000 \\
    \midrule
    GATv2 Dynamic Attention    & \texttimes  & \checkmark  & \checkmark      \\
    SAGEConv (alternating)     & \texttimes  & \texttimes  & \checkmark      \\
    Residual Connections       & \texttimes  & every 2 $l$ & every 2 $l$     \\
    Uncertainty Estimation     & \texttimes  & \texttimes  & \checkmark      \\
    Focal Loss ($\gamma=2.0$)  & \texttimes  & \texttimes  & \checkmark      \\
    Temperature Scaling        & \texttimes  & \checkmark  & \texttimes      \\
    Temporal Smoothing (w=3)   & \texttimes  & \checkmark  & \texttimes      \\
    Graph Readout              & Mean pool   & Node-level  & Attn.\ pool     \\
    Data Augmentation          & \texttimes  & \checkmark  & \checkmark      \\
    Chunk Duration             & 10 s        & 5 s         & 10 s            \\
    \bottomrule
  \end{tabular}
\end{table}

\subsubsection{Training Progression}

The OptimizedRespiratoryGAT stabilised at a weighted F1 of 0.7667, reflecting
the limitations of mean pooling and the absence of class-imbalance
correction. GATv2RespiratoryGAT improved substantially, achieving 0.8788
at epoch 0, demonstrating the benefit of dynamic attention and learned
graph pooling. ImprovedRespiratoryGAT matched this at epoch 0 before
surpassing both prior models at epoch 5, stabilising at a peak
weighted F1 of 0.9384. Training loss fell to 0.0679 at epoch 5 and
reached a minimum of 0.0445 at epoch 11; early stopping triggered at
epoch 20 confirmed the absence of overfitting.

\subsection{Selection of ImprovedRespiratoryGAT}

ImprovedRespiratoryGAT was selected as the final model on five grounds.

\textbf{Superior and stable performance.} The model achieved a peak
validation weighted F1 of 0.9384 which was a 15-point improvement over the
Baseline (0.7667) and a 6-point improvement over GATv2 (0.8788).
The peak was reached at epoch 5 and held across subsequent epochs,
indicating genuine generalisation.

\textbf{Architectural complementarity.} Alternating GATv2Conv and
SAGEConv layers provided dual inductive biases. GATv2 excelled at
fine-grained local pattern recognition, important for wheeze detection,
which exhibited distinct spectral signatures. SAGEConv provided
neighbourhood-level temporal consistency beneficial for crackle
detection, where events spanned multiple consecutive frames.

\textbf{Principled class-imbalance handling.} Focal Loss and
uncertainty-aware weighting jointly addressed the two sources of
imbalance in ICBHI: within-class positive/negative frame imbalance
and between-task imbalance, without requiring manual coefficient tuning.

\textbf{Clinical utility.} The heteroscedastic uncertainty component
produced calibrated confidence scores alongside each prediction,
enabling the model to flag ambiguous cases for clinician review.

\textbf{Training efficiency.} Despite greater architectural complexity,
the model converged within 2 epochs to the F1 target threshold and
exhibited no overfitting, with early stopping at epoch 20 confirming
well-regularised training.

In summary, ImprovedRespiratoryGAT offered the best trade-off across
predictive performance, training stability, class-imbalance robustness,
and clinical interpretability, and constituted the recommended
architecture for respiratory sound event detection in this work.

\subsection{Explainable Artificial Intelligence}

\subsubsection{Attribution Methods}

Two gradient-based attribution methods were applied to all three models
for explainability analysis:

\textbf{Integrated Gradients (IG).} Computes node-wise attributions by
integrating gradients along a linear path from a zero baseline to the input
(Equation~\ref{eq:integrated_gradients}). IG satisfies the
completeness axiom: the sum of attributions equals the difference between
the model output at the input and at the baseline.

\textbf{Grad$\times$Input.} Computes the element-wise product of the
gradient with the input: $\nabla_{\mathbf{x}} F(\mathbf{x}) \odot
\mathbf{x}$. This is computationally cheaper than IG but does not satisfy
the completeness axiom.

For the ImprovedRespiratoryGAT, the learned attention pooling weights
were also extracted as a third attribution source, providing
node-level importance scores before aggregation.

\subsubsection{XAI Evaluation Setup}

Attributions were computed on a single test audio file
(\texttt{107\_2b3\_Ar\_mc\_AKGC417L.wav}) containing both wheeze (33
frames) and crackle (36 frames) annotations across 40 frames. Two
quantitative metrics were computed:

\textbf{Signal-to-Noise Ratio (SNR):} Ratio of mean IG attribution on
ground-truth positive frames to mean attribution on negative frames.
SNR $> 1$ indicates attributions are concentrated on positive frames;
SNR $\approx 1$ indicates uniform distribution.

\textbf{Top-5 IG Overlap:} Fraction of the 5 most attributed frames
that overlap with ground-truth positive frames. Higher overlap indicates
the model is correctly focusing on clinically relevant segments.

\subsubsection{XAI Results}

\begin{table}[!h]
  \renewcommand{\arraystretch}{1.15}
  \caption{XAI Evaluation: Attribution Quality Metrics}
  \label{tab:xai_results}
  \centering
  \begin{tabular}{@{}lcccc@{}}
    \toprule
    \textbf{Model}
      & \textbf{SNR (Wheeze)}
      & \textbf{SNR (Crackle)}
      & \textbf{Top-5 Overlap (Wheeze)}
      & \textbf{Top-5 Overlap (Crackle)} \\
    \midrule
    Baseline GAT   & 0.000 & 0.023 & 0\%  & 40\% \\
    GraphSAGE      & 0.000 & 0.029 & 0\%  & 40\% \\
    Improved GAT   & 1.009 & 1.102 & 60\% & 80\% \\
    \bottomrule
  \end{tabular}
\end{table}

\paragraph{Baseline GAT and GraphSAGE.}
Both node-level models exhibit low IG SNR values ($< 0.03$), indicating
that attributions are approximately uniformly distributed across all
frames rather than concentrated on positive events. The 0\% top-5 overlap
for wheeze suggests these models do not learn to localise wheeze events
in the input features. The 40\% crackle overlap indicates partial
localisation for crackle events, consistent with the higher crackle
ROC-AUC values observed in evaluation.

\paragraph{ImprovedRespiratoryGAT.}
The Improved GAT achieves the highest attribution quality: IG SNR $> 1.0$
for both tasks (wheeze = 1.009, crackle = 1.102), confirming that
attributions are concentrated on ground-truth positive frames. The 60\%
top-5 overlap for wheeze and 80\% for crackle demonstrate strong
localisation ability. This is consistent with the attention pooling
mechanism, which learns to weight diagnostically informative frames more
heavily during graph-level aggregation.

\paragraph{Attention Weights.}
The Improved GAT's attention weights (entropy = 3.686 on a 40-node
graph, maximum possible entropy = $\ln 40 \approx 3.689$) are nearly
uniform, indicating the model distributes attention broadly across the
graph. This near-uniform attention distribution is consistent with the
model's tendency to predict positive for all frames at $t=0.05$---it
does not selectively attend to specific frames but rather averages
across all nodes.

\paragraph{XAI Interpretation.}
The XAI analysis reveals a paradox: the ImprovedRespiratoryGAT produces
the best \textit{attribution quality} (highest SNR and overlap) but the
worst \textit{calibration} (all probabilities below 0.5 at $t=0.5$).
This suggests the model's internal representations correctly distinguish
positive from negative frames, but the attention pooling mechanism
collapses the representation such that all outputs fall in a narrow
low-probability band. The attribution methods succeed in identifying
which frames the model considers important, even though the model's
final probability outputs are poorly calibrated.

\subsection{Hardware Profiling and Edge-Deployment Feasibility}

\begin{table}[!h]
  \renewcommand{\arraystretch}{1.15}
  \caption{Hardware Profiling: Edge-Deployment Feasibility}
  \label{tab:hardware}
  \centering
  \begin{tabular}{@{}lrrrrrr@{}}
    \toprule
    \textbf{Model}
      & \textbf{Params}
      & \textbf{Size (MB)}
      & \textbf{FLOPs (M)}
      & \textbf{Latency (ms)}
      & \textbf{p95 (ms)}
      & \textbf{Throughput} \\
    \midrule
    Baseline GAT   & 463,362 & 1.78 & 11.85 & 6.68  & 22.55 & 149.8 inf/s \\
    GraphSAGE      & 658,434 & 2.53 &  5.92 & 2.71  &  3.75 & 368.7 inf/s \\
    Improved GAT   & 677,765 & 2.61 & 13.50 & 11.13 & 31.02 &  89.8 inf/s \\
    \bottomrule
  \end{tabular}
\end{table}

\noindent All three GNN models require fewer than 3\,MB for checkpoint
storage and negligible runtime memory during inference. Including the
Wav2Vec2 feature extractor ($\sim$350\,MB), the total system memory
requirement is approximately 400\,MB---feasible on mid-range Android devices
(2020+) with 4\,GB RAM. The GNN models themselves require 5.9--13.5\,MFLOPs
per inference, negligible compared to the Wav2Vec2 encoder ($\sim$2.5\,GFLOPs
per 10-second chunk). On CPU, all three GNN models achieve sub-12\,ms median
inference per chunk, confirming that the graph reasoning component is not
the latency bottleneck. At 90--369 inferences per second, all models process
audio chunks faster than real-time, confirming feasibility for continuous
ward monitoring.

\subsection{Data Augmentation}

Data augmentation was applied during training for the GraphSAGE and
ImprovedRespiratoryGAT models using the \texttt{AudioAugmentation} class
(\texttt{model\_comparisons/gnn\_augmentation.py}). The Baseline GAT was
trained without augmentation. The augmentation pipeline applies a random
transformation with probability $p=0.3$ per sample, selecting uniformly from:
time stretching (rate $\in [0.85, 1.15]$), pitch shifting (n\_steps $\in
\{-2, \ldots, +2\}$), additive Gaussian noise (level $\in [0, 0.05]$),
time masking (10--30\% of frames zeroed), frequency masking (10--30\% of
frequency bins zeroed), or no augmentation (10\% probability). Raw audio
augmentations are applied before Wav2Vec2 feature extraction; embedding
augmentations are applied directly to the extracted embeddings.